\chapter{Lý thuyết tiếp cận}
\newpage
\section{Phương pháp tính động lực học}
\paragraph{}Có nhiều phương pháp dùng để tính động lực học. Ví dụ như: phương pháp Newton, phương pháp Lagrange, phương pháp theo năng lượng,... Nhưng trong khuôn khổ bài đề tài này, phương pháp Newton được sử dụng với những ưu điểm có thể kể đến như sau:
\begin{itemize}
	\item Phương pháp Newton sử dụng các phương pháp tính cơ học thông thường.
	\item Các công thức và hệ phương trình trong quá trình tính không quá phức tạp.
	\item Kết quả tính động lực học của mô hình con lắc ngược phổ biến hiện nay trong các tài liệu tham khảo được sử dụng để kiểm tra sai sót trong quá trình tính toán động lực học của mô hình xe hai bánh tự cân bằng.
\end{itemize}
\paragraph{}Bên cạnh ba ưu điểm trên, phương pháp Newton vẫn có nhược điểm là phải tuyến tính hoá tính toán tại vị trí góc $\theta$ = 0${^\circ}$. Tuy nhiên việc này không gây ảnh hưởng quá nhiều tới mô hình xe hai bánh cân bằng được xây dựng trong đề tài này, vì mô hình chỉ hoạt động xung quanh vị trí 0${^\circ}$ ($\pm$10${^\circ}$).
\subsection{Nền tảng lý thuyết từ con lắc ngược}
\begin{center}
	\begin{figure}[htp]
		\begin{center}
			\includegraphics[scale=.7]{image/hinh5.png}
		\caption{Mô hình con lắc ngược}
				\end{center}
	\end{figure}
\end{center}
\paragraph{}Ta xem xét mô hình toán học của con lắc ngược thông qua một hệ thống bao gồm một con lắc ngược được gắn vào một xe đẩy có gắn động cơ. Hệ thống con lắc ngược cũng tương tự như xe máy đứng: gia tốc cao sẽ ảnh hướng đến góc nghiêng của xe. Con lắc ngược là một hệ thống con lắc mà khối tâm nằm ngay trên trục của nó, cho nên sự thăng bằng có thể đạt được là không hề ổn định, rất khó đạt đến mức cân bằng. Hệ thống con lắc ngược là một ví dụ thường thấy trong các giáo trình về hệ thống điều khiển và các tài liệu nghiên cứu khác. Sự phổ biến của nó xuất phát một phần từ thực tế là nó không ổn định nếu không có sự kiểm soát, nghĩa là, con lắc đơn giản sẽ rơi xuống nếu như xe đẩy không được di chuyển để cân bằng nó. Ngoài ra, các phương trình động lực học của hệ là phi tuyến. Mục tiêu của hệ thống điều khiển là cân bằng được  con lắc ngược bằng cách tác dụng một lực đẩy lên xe đẩy mà con lắc được gắn vào. Hiện nay có nhiều thuật toán điều khiển được sử dụng để có thể giải quyết được vấn đề này, chẳng hạn như: bộ điều khiển PID, mạng neutral, điều khiển mờ, thuật toán di truyền,... 
\paragraph{}
Mô hình nghiên cứu gồm hai phần: 
\begin{enumerate}
	\item Một con lắc gắn với xe bởi một khớp bản lề.
	\item Lực F tác động vào xe theo phương nằm ngang. 
\end{enumerate}
\paragraph{}Mô hình gồm hai biến đầu ra:
\begin{enumerate}
	\item Độ dịch chuyển của xe (x).
	\item Góc lệch của con lắc so với phương thẳng đứng ($\theta$).
\end{enumerate}
	\begin{figure}[htp]
		\begin{center}
			\includegraphics[scale=.8]{image/hinh6.png}
			\caption{Sơ đồ khối các biến lối vào và biến lối ra của mô hình con lắc ngược}
		\end{center}
	\end{figure}
\paragraph{}Để giải quyết bài toán này, trước hết ta có các tham số như sau: 
\begin{table}[H]
	\begin{center}
		\begin{tabular}{lll}
			M &  & Khối lượng xe (kg)                       \\
			m &  & Khối lượng con lắc (kg)                  \\
			b &  & Ma sát của xe (N)                        \\
			L &  & Chiều dài 1/2 con lắc (m)                \\
			I &  & Momen quán tính của con lắc (Nm)         \\
			F &  & Lực tác động vào xe (N)                  \\
			x &  & Vị trí của xe (m)                        \\
			$\theta$ &  & Góc  của con lắc so với phương thẳng đứng (rad)
		\end{tabular}
	\end{center}
\end{table}
\paragraph{}Phân tích các lực tác dụng lên xe và con lắc, ta có:
\begin{figure}[H]
	\begin{center}
		\includegraphics[scale=.6]{image/hinh7.png}
		\caption{Sơ đồ các lực tác dụng lên xe và con lắc}
	\end{center}
\end{figure}
\paragraph{}Tổng hợp lực tác dụng lên xe theo phương ngang:
\begin{equation}
M\ddot{x} + b\dot{x} + N = F
\end{equation}
\paragraph{}Tổng hợp lực tác dụng lên con lắc theo phương ngang: 
\begin{equation}
N = m\ddot{x} + mL\ddot{\theta}\cos{\theta} - mL\dot{\theta}^{2}\sin{\theta}
\end{equation}
\paragraph{}Thay (2.2) vào (2.1) ta được phương trình chuyển động đầu tiên của hệ:
\begin{equation}
(M + m)\ddot{x} + b\dot{x} + mL\ddot{\theta}\cos{\theta} - mL\dot{\theta}^{2}\sin{\theta} = F
\end{equation}
\paragraph{}Tổng hợp lực tác dụng lên con lắc theo phương vuông góc:
\begin{equation}
P\sin{\theta} + N\cos{\theta} + mg\sin{\theta} = mL\ddot{\theta} + m\ddot{x}\cos{\theta}
\end{equation}
\paragraph{}Tổng momen tại khối tâm con lắc: 
\begin{equation}
-PL\sin{\theta} - NL\cos{\theta} = I\ddot{\theta}
\end{equation}
\paragraph{}Kết hợp phương trình (2.4) và (2.5) ta được phương trình chuyển động thứ hai của hệ: 
\begin{equation}
(I + mL^{2})\ddot{\theta} - mgL\sin{\theta} = - mL\ddot{x}\cos{\theta}
\end{equation}
\paragraph{}Từ (2.3) và (2.6) ta được hệ phương trình chuyển động của hệ:
\begin{equation}
\begin{cases}
(M + m)\ddot{x} + b\dot{x} + mL\ddot{\theta}\cos{\theta} - mL\dot{\theta}^{2}\sin{\theta} = F \\
(I + mL^{2})\ddot{\theta} - mgL\sin{\theta} = - mL\ddot{x}\cos{\theta}
\end{cases}
\end{equation}
\paragraph{}Xấp xỉ tuyến tính hoá hai phương trình của hệ (2.7) tại 0${^\circ}$:
\begin{equation}
\begin{cases}
(M + m)\ddot{x} + b\dot{x} + mL\ddot{\theta} = F  \\
(I + mL^{2})\ddot{\theta} - mgL\theta = - mL\ddot{x}
\end{cases}
\end{equation}
\paragraph{}Biểu diễn hệ phương trình chuyển động của hệ sau khi đã được tuyến tính hoá dưới dạng hàm trạng thái: 
\begin{align}
\begin{bmatrix}
\dot{x}\\
\ddot{x}\\
\dot{\theta}\\
\ddot{\theta}
\end{bmatrix} 
&=
\begin{bmatrix}
0 & 1 & 0 & 0\\
0 & \dfrac{-(I + mL^{2})b}{I(M + m) + MmL^{2}} & \dfrac{-m^{2}gL^{2}}{I(M + m) + MmL^{2}} & 0\\
0 & 0 & 0 & 1 \\
0 & \dfrac{mLb}{I(M + m) + MmL^{2}} & \dfrac{mgL(M + m)}{I(M + m) + MmL^{2}} & 0
\end{bmatrix} 
\begin{bmatrix}
x\\
\dot{x}\\
\theta\\
\dot{\theta}
\end{bmatrix}\\
\nonumber
&+
\begin{bmatrix}
0\\
\dfrac{I + mL^{2}}{I(M + m) + MmL^{2}} \\
0\\
\dfrac{-mL}{I(M + m) + MmL^{2}}
\end{bmatrix}F
\end{align}
\paragraph{}Biến đầu ra:
\begin{align}
\begin{bmatrix}
x\\
\theta
\end{bmatrix}
=
\begin{bmatrix}
1 & 0 & 0 & 0\\
0 & 0 & 1 & 0 
\end{bmatrix}
\begin{bmatrix}
x\\
\dot{x}\\
\theta\\
\dot{\theta}
\end{bmatrix}
+
\begin{bmatrix}
0\\
0
\end{bmatrix}F
\end{align}
\paragraph{}Nếu cho b $\ll$ 1 và I $\ll$ 1, ta được dạng đơn giản hơn của hàm trạng thái: 
\begin{align}
\begin{bmatrix}
\dot{x}\\
\ddot{x}\\
\dot{\theta}\\
\ddot{\theta}
\end{bmatrix}
=
\begin{bmatrix}
0 & 1 & 0 & 0\\
0 & 0 & \dfrac{-mg}{M} & 0\\
0 & 0 & 0 & 1\\
0 & 0 & \dfrac{g(M + m)}{ML}
\end{bmatrix}
\begin{bmatrix}
x\\
\dot{x}\\
\theta\\
\dot{\theta}
\end{bmatrix}
+
\begin{bmatrix}
0\\
\dfrac{1}{M}\\
0\\
\dfrac{-1}{ML}
\end{bmatrix}F
\end{align}
\subsection{Động lực học mô hình xe tự cân bằng}
\paragraph{}Mô hình xe tự cân bằng bao gồm một khung gầm mang theo các thành phần điện tử như động cơ DC, bảng mạch, các cảm biến...; một thanh thép dọc gắn với thân xe, trên thanh dọc có lắp một số quả nặng để "mô phỏng" khối lượng của người lái; các bánh xe được gắn trực tiếp với động cơ DC.
\begin{center}
	\begin{figure}[htp]
		\begin{center}
			\includegraphics[scale=.7]{image/hinh14.png}
			\caption{Định nghĩa các biến đầu vào và nhiễu của hàm trạng thái}
		\end{center}
	\end{figure}
\end{center}
\paragraph{}Hình 2.4 cho thấy chiếc xe có 3 bậc tự do. Nó có thể xoay quanh trục z (pitch), chuyển động này được mô tả bằng góc $\theta_{P}$ với vận tốc góc tương ứng là $\omega_{P}$. Chuyển động tuyến tính của khung gầm được đặc trưng bởi vị trí $x_{RM}$ và vận tốc $v_{RM}$. Ngoài ra, chiếc xe còn có thể xoay quanh trục thẳng đứng y (yaw) với sự kết hợp của góc $\delta$ và vận tốc góc $\dot{\delta}$. Sáu biến của hàm trạng thái này mô tả đầy đủ động lực học của một hệ 3 bậc tự do. Tuy nhiên chúng ta chỉ xem xét chuyển động xoay quanh trục z của xe nên chỉ cần quan tâm đến đạo hàm bậc nhất và bậc hai của $x_{RM}$ và $\theta_{P}$ để mô tả động lực học của hệ 2 bậc tự do. 
\paragraph{}Chiếc xe được điều khiển bằng cách áp dụng momen xoắn $C_{L}$ và $C_{R}$ đến các bánh xe tương ứng.
\paragraph{}Để có thể kiểm soát được hệ thống thành công, các biến của hàm không gian phải được xác định, hoặc thông qua đo lường trực tiếp hoặc phải được quan sát cẩn thận. Góc và tốc độ góc của chuyển động theo trục z có thể được xác định dễ dàng bằng cảm biến tích hợp gia tốc kế và con quay hồi chuyển. 
\paragraph{}Hệ thống điều khiển dựa trên bộ điều khiển trạng thái nhằm điểu khiển sự ổn định xung quanh trục z nằm ngang (pitch). Bộ điều khiển sẽ tạo ra một giá trị riêng của momen xoắn cho từng bánh xe, tín hiệu này sẽ được đưa đến động cơ tương ứng. 
\begin{center}
	\begin{figure}[htp]
		\begin{center}
			\includegraphics[scale=.7]{image/hinh15.png}
			\caption{Giản đồ lực của mô hình xe hai bánh tự cân bằng}
		\end{center}
	\end{figure}
\end{center}
\paragraph{}Dựa vào giản đồ lực hình 2.5, ta có:
\begin{itemize}
	\item Đối với bánh trái (tương tự với bánh phải): \\
\begin{equation}
	\ddot{x_{RL}}M_{RL} = H_{TL} - H_{L}
\end{equation}
\begin{equation}
\ddot{y_{RL}}M_{RL} = V_{TL} - M_{RL}g - V_{L}
\end{equation}
\begin{equation}
\ddot{\theta_{RL}}J_{RL} = C_{L} - H_{TL}R
\end{equation} 
\item Đối với thân xe: \\
\begin{equation}
\ddot{x_{P}}M_{P} = H_{R} + H_{L} 
\end{equation}
\begin{equation}
\ddot{y_{P}}M_{P} = V_{L} + V_{R} - M_{P}g + F_{C\theta}
\end{equation}
\begin{equation}
J_{P}\ddot{\theta} = (V_{L} + V_{R})L\sin{\theta} - (H_{L} + H_{R})L\cos{\theta} - (C_{L} + C_{R})
\end{equation}\\
Trong đó: \\

\begin{table}[H]
	\begin{center}
	\begin{tabular}{lll}
		J$_{P}$       & : Momen quán tính của khung gầm xe (Nm$^{2}$)                   \\
		J$_{RL}$, J$_{RR}$       & : Momen quán tính của bánh xe trái và bánh xe phải (Nm$^{2}$)                             \\
		M$_{P}$        & : Khối lượng của khung gầm xe (kg)                           \\
		M$_{RL}$, M$_{RR}$          & : Khối lượng của bánh xe trái và bánh xe phải (kg)                                \\
		J$_{W}$ & : Momen trung bình của bánh (Nm$^{2}$)  \\
		M$_{W}$ & : Khối lượng trung bình của bánh (kg)\\
		R          & : Bán kính của bánh xe (m)                              \\
		D          & : Khoảng cách giữa hai bánh xe (m)                         \\
		L          & : Khoảng cách giữa trục z và trọng tâm của khung gầm xe (m) \\
		$\theta$   & : Góc lật (rad)        \\
		$\theta_{W}$ & : Góc lật trung bình (rad)                                     
	\end{tabular}
\end{center}
\end{table}
\paragraph{}Tổng quá ta được phương trình: 
\begin{equation}
J_{P}\ddot{\theta} = (V_{L} + V_{R})L\sin{\theta} - (H_{L} + H_{R})L\cos{\theta} - (C_{L} + C_{R})
\end{equation} 
\begin{equation}
x_{P}M_{P} = H_{L} + H_{R}
\end{equation}
\begin{equation}
\ddot{y_{P}}M_{P} = V_{L} + V_{R} - M_{P}g + \dfrac{C_{L}+C_{R}}{L}\sin{\theta}
\end{equation}
\end{itemize}
\paragraph{} Thế (2.19) và (2.20) vào (2.18) ta có: 
\begin{equation}
J_{P}\ddot{\theta} = (\ddot{y_{P}}M_{P} + M_{P}g - \dfrac{C_{L}+C_{R}}{L})L\sin^{2}{\theta} - \ddot{x_{P}}M_{P}L\cos{\theta} - (C_{L} + C_{R})
\end{equation}
\begin{equation}
J_{P}\ddot{\theta} = M_{P}L(\ddot{y_{P}}\sin{\theta} - \ddot{x_{P}}\cos{\theta}) + M_{P}gL\sin{\theta} - (1 + \sin^{2}{\theta})(C_{L}+C_{R})
\end{equation}
\paragraph{}Xét với mỗi bánh trái và bánh phải: 
\begin{equation}
\begin{cases}
\ddot{x_{RL}}M_{RL} = H_{TL} - H_{L} \\
\ddot{x_{RR}}M_{RR} = H_{TR} - H_{R} 
\end{cases}
\end{equation}
\paragraph{}Mà: 
\begin{equation}
M_{RL} = M_{RR} = \dfrac{M_{W}}{2} ; \ddot{x_{RL}} + \ddot{x_{RR}} = 2\ddot{x_{RM}}
\end{equation}
\paragraph{}Cộng hai vế của hệ phương trình (2.23) ta được:
\begin{equation}
2M_{W}\ddot{x_{RM}} = H_{TL} + H_{TR} - H_{L} - H_{R}
\end{equation}
\paragraph{}Lại có:
\begin{equation}
J_{RL}\ddot{\theta_{RL}} = C_{L} - H_{TL}R ; \ddot{x_{P}}M_{P} = H_{L} + H_{R}
\end{equation}
\paragraph{}Thay vào (2.24) ta được:
\begin{equation}
2M_{W}\ddot{x_{RM}} = \dfrac{C_{L}+C_{R} - (J_{RL}\ddot{\theta_{RL}} + J_{RR}\ddot{\theta_{RR}})}{R} - \ddot{x_{P}}M_{P}
\end{equation}
\paragraph{}Đối với momen xoắn của bánh trái và bánh phải ta có:
\begin{equation}
\begin{cases}
C_{L} = \frac{1}{2}C_{\theta} + \frac{1}{2i}C_{\delta}\\
C_{R} = \frac{1}{2}C_{\theta} - \frac{1}{2i}C_{\delta}
\end{cases}
\end{equation}
\paragraph{}Cộng hai vế của hệ phương trình (2.28) ta được C$_{\theta}$ là giá trị trung bình của C$_{L}$ và C$_{R}$:
\begin{equation}
C_{L} + C_{R} = C_{\theta}
\end{equation}
\paragraph{}Phương trình (2.27) trở thành:
\begin{equation}
2M_{W}\ddot{x_{RM}} = \dfrac{C_{\theta}}{R} - 2\dfrac{J_{W}\ddot{\theta_{W}}}{R} - \ddot{x_{P}}M_{P}
\end{equation}
\paragraph{}Có:
\begin{equation}
\begin{cases}
y_{P} = -L(1-\cos{\theta})\\
x_{P} = x_{RM} + L\sin{\theta}
\end{cases}
\end{equation}
\begin{equation}
\begin{cases}
\dot{y_{P}} = -\theta L\sin{\theta}\\
\dot{x_{P}} = \dot{x_{RM}} + \theta L\cos{\theta}
\end{cases}
\end{equation}
\begin{equation}
\begin{cases}
\ddot{y_{P}} = -\ddot{\theta} L\sin{\theta} - \dot{\theta^{2}} L\cos{\theta}\\
\ddot{x_{P}} = \ddot{x_{RM}} + \ddot{\theta} L\cos{\theta} - \dot{\theta^{2}} L\sin{\theta}
\end{cases}
\end{equation}
\paragraph{}Do đó:
\begin{equation}
\ddot{y_{P}}\sin{\theta} - \ddot{x_{P}}\cos{\theta} = -\ddot{\theta}L - \ddot{x_{RM}}\cos{\theta}
\end{equation}
\paragraph{}Thay (2.33) vào (2.30) và (2.34) vào (2.22) ta được hệ phương trình:
\begin{equation}
\begin{cases}
J_{P}\ddot{\theta} = -M_{P}L^{2}\ddot{\theta} + M_{P}gL\sin{\theta} - (1 + \sin{\theta})C_{\theta} -\ddot{x_{RM}}LM_{P}\cos{\theta}\\
2M_{W}\ddot{x_{RM}} = \dot{\theta^{2}}M_{P}L\sin{\theta} - \ddot{\theta}M_{P}L\cos{\theta} - 2\dfrac{J_{W}\ddot{\theta_{W}}}{R} + \dfrac{C_{\theta}}{R} - M_{P}\ddot{x_{RM}}
\end{cases}
\end{equation}
\paragraph{}Momen quán tính của thân xem như là một đoạn thẳng hình trụ có chiều dài 2L, bánh xe xem như là một đĩa tròn xoay:
\begin{equation}
\begin{cases}
J_{P} = \dfrac{M_{P}L^{2}}{3}\\
J_{W} = \dfrac{M_{W}R^{2}}{2}
\end{cases}
\end{equation}
\paragraph{}Lúc này hệ phương trình (2.35) trở thành:
\begin{equation}
\begin{cases}
\dfrac{M_{P}L^{2}}{3}\ddot{\theta} = -M_{P}L^{2}\ddot{\theta} + M_{P}gL\sin{\theta} - (1 + \sin{\theta})C_{\theta} -\ddot{x_{RM}}LM_{P}\cos{\theta}\\
(2M_{W}+M_{P})\ddot{x_{RM}} =  \dot{\theta^{2}}M_{P}L\sin{\theta} - \ddot{\theta}M_{P}L\cos{\theta} - 2\dfrac{\dfrac{M_{W}R^{2}\ddot{\theta}}{2}}{R} + \dfrac{C_{\theta}}{R}
\end{cases}
\end{equation}
\begin{equation}
\begin{cases}
\dfrac{4M_{P}L^{2}}{3}\ddot{\theta} = + M_{P}gL\sin{\theta} - (1 + \sin{\theta})C_{\theta} -\ddot{x_{RM}}LM_{P}\cos{\theta}\\
(2M_{W}+M_{P})\ddot{x_{RM}} =  \dot{\theta^{2}}M_{P}L\sin{\theta} - \ddot{\theta}M_{P}L\cos{\theta} - M_{W}R\ddot{\theta} + \dfrac{C_{\theta}}{R}
\end{cases}
\end{equation}
\paragraph{}Xấp xỉ các giá trị lượng giác ta được: 
\begin{equation}
\begin{cases}
\dfrac{4M_{P}L\ddot{\theta}}{3} - M_{P}g\theta + \dfrac{C_{\theta}}{L} + M_{P}\ddot{x_{RM}} = 0\\
(2M_{W}+M_{P})\ddot{x_{RM}} + (M_{P}L+M_{W}R)\ddot{\theta} = \dfrac{C_{\theta}}{R}
\end{cases}
\end{equation}
\begin{equation}
\begin{cases}
\dfrac{4L\ddot{\theta}}{3} - g\theta + \dfrac{C_{\theta}}{M_{P}L} + \ddot{x_{RM}} = 0\\
M_{P}\ddot{x_{RM}} + \dfrac{M_{P}}{2M_{W}+M_{P}}(M_{P}L+M_{W}R)\ddot{\theta} = \dfrac{M_{P}}{2M_{W}+M_{P}}\dfrac{C_{\theta}}{R}
\end{cases}
\end{equation}
\paragraph{}Giải hệ phương trình ta được: 

\begin{subnumcases}
\ddot{x_{RM}} = -\dfrac{4L\ddot{\theta}}{3} + g\theta - \dfrac{C_{\theta}}{M_{P}L}\\
\left[\dfrac{4M_{P}L}{3}-\dfrac{M_{P}}{2M_{W}+M_{P}}(M_{P}L+M_{W}R)\ddot{\theta}-M_{P}g\theta\right] = -\dfrac{M_{P}}{2M_{W}+M_{P}}\dfrac{C_{\theta}}{R}-\dfrac{C_{\theta}}{L}
\end{subnumcases}

\paragraph{}Đặt: 
\begin{equation}
X = \dfrac{4M_{P}L}{3}-\dfrac{M_{P}}{2M_{W}+M_{P}}(M_{P}L+M_{W}R)\ddot{\theta}
\end{equation}
\begin{equation}
Y = \dfrac{M_{P}}{2M_{W}+M_{P}}\dfrac{1}{R}+\dfrac{1}{L}
\end{equation}
\paragraph{}Thay (2.42) và (2.43) vào (2.41b) ta được: 
\begin{equation}
\ddot{\theta} = \dfrac{M_{P}g}{X}\theta - \dfrac{YC_{\theta}}{X}
\end{equation}
\paragraph{}Thay (2.44) vào (2.41a) ta được: 
\begin{equation}
\ddot{x_{RM}} = \left(\dfrac{-4LM_{P}g}{3X}+g\right)\theta + \left(\dfrac{4LY}{3X}-\dfrac{1}{M_{P}L}\right)C_{\theta}
\end{equation} 
\paragraph{}Từ đó ta có hệ phương trình biến trạng thái: 
\begin{equation}
\begin{cases}
\ddot{\theta} = \dfrac{M_{P}g}{X}\theta - \dfrac{YC_{\theta}}{X}\\
\ddot{x_{RM}} = \left(\dfrac{-4LM_{P}g}{3X}+g\right)\theta + \left(\dfrac{4LY}{3X}-\dfrac{1}{M_{P}L}\right)C_{\theta}
\end{cases}
\end{equation}
\paragraph{}Hàm trạng thái:
\begin{align}
\begin{bmatrix}
\dot{x_{W}}\\
\ddot{x_{W}}\\
\dot{\theta}\\
\ddot{\theta}
\end{bmatrix}
=
\begin{bmatrix}
0&1&0&0\\
0&0&g\left(1-\dfrac{4LM_{P}}{3X}\right)&0\\
0&0&0&1\\
0&0&\dfrac{M_{P}g}{X}&0
\end{bmatrix}
\begin{bmatrix}
x_{W}\\
\dot{x_{W}}\\
\theta\\
\dot{\theta}
\end{bmatrix}
+
\begin{bmatrix}
0\\
\left(\dfrac{4LY}{3X}-\dfrac{1}{M_{P}L}\right)\\
0\\
-\dfrac{Y}{X}
\end{bmatrix}
\end{align}
\begin{align}
\begin{bmatrix}
x_{W}\\
\theta
\end{bmatrix}
=
\begin{bmatrix}
1&0&0&0\\
0&0&1&0
\end{bmatrix}
\begin{bmatrix}
x_{W}\\
\dot{x_{W}}\\
\theta\\
\dot{\theta}
\end{bmatrix}
+
\begin{bmatrix}
0\\
0
\end{bmatrix}
C_{\theta}
\end{align}
\section{Kỹ thuật lý thuyết điều khiển hiện đại}
\subsection{Thiết kế cổ điển và hiện đại}
\paragraph{}Một nhà thiết kế thời trang điều chỉnh trang phục để đáp ứng thị hiếu của những người tiêu dùng sành điệu, nắm rõ trong đầu loại quần áo này sẽ được mặc vào mùa nào, loại quần áo kia là để dành cho dịp gì. Tương tự như vậy, một kỹ sư điều khiển hệ thống thiết kế hệ thống điều khiển để đáp ứng các mục tiêu mong muốn, ghi nhớ các vấn đề như "ở đâu" và "làm thế nào" mà hệ thống điều khiển thực hiện. Chúng ta cần một hệ thống điều khiển bởi vì chúng ta không thích cách mà máy móc đang vận hành, và bằng cách thiết kế một hệ thống điều khiển chúng ta có thể thử thay đổi cách vận hành của chúng sao cho phù hợp với yêu cầu của mình. Thiết kế đề cập đến quá trình thay đổi các tham số của hệ thống điều khiển để đáp ứng các mục tiêu về độ ổn định, hiệu suất và độ bền. Các tham số thiết kế có thể là các hằng số chưa biết trong hàm truyền của bộ điều khiển hoặc biểu diễn không gian trạng thái của nó. 
\paragraph{}Có hai bộ phận chính trong thiết kế hệ thống điều khiển, cụ thể là, cổ điển và hiện đại, tác động trực tiếp đến các ứng dụng điều khiển tự động. 
\paragraph{}Phương pháp thiết kế cổ điển bao gồm cách thay đổi hàm truyền của bộ điểu khiển cho đến khi đạt được hiệu suất vòng kín mong muốn. Những dấu hiệu nhận biết cổ điển của hiệu suất vòng kín là khả nẳng đáp ứng tần số , hoặc vị trí của các cực của vòng kín. Đối với hệ thống điều khiển bậc cao, bằng cách thay đổi một số lượng giới hạn hằng số trong hàm truyền của bộ điều khiển, có thể làm thay đổi phần lớn các vị trí cực vòng kín, nhưng không phải là tất cả. Đây chính là một hạn chế lớn của phương pháp thiết kế cổ điển. 
\paragraph{}Đối với phương pháp thiết kế hiện đại, tham số thiết kế được thay đổi là biểu diễn không gian trạng thái của bộ điều khiển. Nguyên tắc kiểm tra hệ thống ổn định bằng hàm trạng thái: tìm nghiệm $\lambda$ của phương trình $\left|\lambda I - A\right| = 0$, với A là một thành phần của hàm trạng thái. Nếu một trong những giá trị riêng của A là $\lambda > 0$ thì vòng kín  của hệ điều khiển không ổn định.
\paragraph{}Khả năng điều khiển của hệ thống có thể được xác lập khi hệ thống có được bất kỳ trạng thái ban đầu nào, $x(t_{0})$, đến trạng thái cuối bất kỳ nào, $x(t_{f})$, trong một thời gian xác định, $(t_{f} - t_{0})$, với mỗi một giá trị của véc-tơ ngõ vào, $u(t)$, $t_{0} < t < t_{f}$. Điều này cực kỳ quan trọng, vì có thể thay đổi được hệ thống không điều khiển được từ trạng thái ban đầu đến trạng thái cuối cùng, hay lấy một lượng thời gian không xác định làm thay đổi hệ thống không điều khiển được bằng cách dùng véc-tơ ngõ vào $u(t)$. Khả năng điều khiển được của hệ thống có thể dễ dàng kiểm tra nếu tách riêng ra các phương trình trạng thái của hệ thống. Mỗi phương trình trạng thái vô hướng tách riêng đều tương đương với một hệ thống phụ. 
\subsection{Định lý về khả năng điều khiển}
\begin{dl}Một hệ thống tuyến tính, bất biến theo thời gian được mô tả bằng phương trình trạng thái ma trận, \textbf{x$^{(1)}$(t) = Ax(t) + Bu(t)} điều khiển được nếu và chỉ nếu ma trận kiểm tra khả năng điều khiển \textbf{P = [AB; A$^{2}$B; A$^{3}$B; ...; A$^{n-1}$B]} có hạng là n, bằng với bậc của hệ thống. 
\end{dl}
\paragraph{}Hạng của ma trận P, được định nghĩa là cấp cao nhất của các định thức con khác 0 có trong P. Nếu P là ma trận vuông, định thức lớn nhất hình thành P là |P|. Nếu P không phải là ma trận vuông, thì định thức lớn nhất hình thành P được tính bằng bằng cách cho tất cả số hàng bằng với số cột hoặc tất cả số cột bằng với số hàng của P, từ đó ta sẽ có được hạng của ma trận P. Chú ý đối với hệ thống bậc \textit{n} với \textit{r} lối vào thì kích thước của ma trận kiểm tra khả năng điều khiển P là \textit{n}$\times$\textit{m}. Định thức con khác 0 lớn nhất của P có thể là thứ nguyên của \textit{n}. Do đó, hạng của P có thể nhỏ hơn hoặc bằng với \textit{n}. 
\paragraph{}Một minh chứng rõ ràng của định lý kiểm tra khả năng điều khiển đại số có thể thấy trong định lý Friedland (phần đại số ma trận). Một dạng tương tự khác của định lý có thể áp dụng cho hệ thống thay đổi theo thời gian. Từ đó có thể viết lại ma trận kiểm tra khả năng điều khiển thay đổi theo thời gian như sau: \textbf{P(t) = [B(t); A(t)B(t); A$^{2}$(t)B(t); A$^{3}$(t)B(t); ...; A$^{n-1}$(t)B(t)]} và kiểm tra hạng của P đối với mọi thời điểm, $t>t_{0}$, cho hệ thống tuyến tính biến đổi theo thời gian. Nếu tại bất kỳ thời điểm t nào đó, hạng của P(t) nhỏ hơn \textit{n}, hệ thống sẽ không thể điều khiển được. Tuy nhiên, việc dùng ma trận kiểm tra khả năng điều khiển biến đổi theo thời gian vừa nêu cần phải chú ý đến hệ số trạng thái biến đổi nhanh theo thời gian, bởi vì việc kiểm tra có thể thực thi tại các bước thời gian rời rạc - tốt hơn so với thực thi trên toàn bộ thời gian liên tục, và trong vài khoảng thời gian (ngắn hơn bước thời gian) hệ thống có lẽ sẽ mất kiểm soát, không điều khiển được. 
\paragraph{}Tóm lại, định lý về khả năng điều khiển có thể được hiểu như sau:
\begin{center}
	Nếu |P| $\geq$ rank(P) thì hệ thống điều khiển được
\end{center}
\paragraph{}Việc kiểm tra khả năng điều khiển gồm nhận biết hạng của ma trận P, kiểm tra xem liệu nó có bằng n hay không. Điều này yêu cầu phải thực hiện một quá trình tính toán phức tạp và dài dòng nếu như tự tính bằng tay. Tuy nhiên, phần mềm MATLAB cung cấp sẵn lệnh \textit{rank(P)} giúp cho việc nhận biết hạng của ma trận dễ dàng hơn. Hơn nữa, công cụ Control System Toolbox cho phép hình thành trực tiếp ma trận kiểm tra khả năng điều khiển P bằng cách dùng lệnh \textit{ctrb}. Ví dụ: \textit{P = ctrb(A,B)} hoặc \textit{P = ctrb(sys)}, trong đó A và B là ma trận hệ số trạng thái của hệ thống mà hàm truyền của nó là \textit{sys}.
\paragraph{}Nguyên nhân khiến cho hệ thống không điều khiển được có thể là do tính chính xác, chẳng hạn cách dùng các giá trị trạng thái thừa (nghĩa là các giá trị trạng thái nhiều hơn bậc của hệ thống). Giá trị trạng thái thừa không ảnh hưởng trực tiếp bởi ngõ vào hệ thống, sẽ không tạo ra không gian trạng thái để điều khiển, cho dù hệ thống có thể điều khiển được về mặt vật lý. Nguyên nhân đôi khi làm mất khả năng điều khiển được là có quá nhiều cực đối xứng trong mô hình toán học của hệ thống. 
\subsection{Thiết kế gán cực hồi tiếp biến toàn trạng thái} 
\paragraph{}Hệ thống điều khiển vòng kín

 được thiết kế bởi bộ điều khiển để đặt các cực tại vị trí mong muốn sẽ làm thay đổi đặc tính của hệ điều khiển. Phương pháp thiết kế cổ điển dùng hàm chuyển giao bộ điều khiển với một vài tham số thiết kế không đủ đặt vào tất cả các cực của vòng kín tại vị trí mong muốn. Phương pháp hàm trạng thái dùng thông tin hồi tiếp toàn trạng thái cung cấp đủ tham số thiết kế bộ điều khiển để di chuyển các cực vòng kín độc lập với nhau. Thông tin hồi tiếp toàn trạng thái được phát ra từ véc-tơ ngõ vào của bộ điều khiển, $u(t)$, tuân theo định luật điều khiển như sau: 
\begin{equation}
u(t) = K[x_{d}(t) - x(t)] - K_{d}x_{d}(t) - K_{n}x_{n}(t)
\end{equation}
\paragraph{}Trong đó, $x(t)$ là véc-tơ trạng thái của hệ thống, $x_{d}(t)$ là véc-tơ trạng thái mong muốn, $x_{n}(t)$ là véc-tơ trạng thái nhiễu và $K, K_{d}, K_{n}$ là ma trận tăng cường của bộ điều khiển. Véc-tơ trạng thái mong muốn, $x_{d}(t)$, và véc-tơ trạng thái nhiễu, $x_{n}(t)$, được tạo ra từ tiến trình ngoài và hoạt động như các ngõ vào hệ thống điều khiển. Nhiệm vụ của bộ điều khiển là thu được véc-tơ trạng thái mong muốn ở tình trạng ổn định, trong khi đó vẫn phản ứng chống lại các ảnh hưởng do nhiễu gây ra. Véc-tơ ngõ vào, $u(t)$, ứng dụng cho hệ thống được mô tả bởi phương trình ngõ ra và trạng thái sau: 
\begin{equation}
x^{(1)}(t) = Ax(t) + Bu(t) + Fx_{n}(t)
\end{equation}
\begin{equation}
y(t) = Cx(t) + Du(t) + Ex_{n}(t)
\end{equation}
\paragraph{}Trong đó F và E là ma trận hệ số nhiễu trong phương trình ngõ ra và trạng thái. Thiết kế hệ thống điều khiển dùng thông tin hồi tiếp toàn trạng thái đòi hỏi hệ điều khiển được mô tả bằng phương trình (2.50) phải có khả năng điều khiển, mặt khác ngõ vào điều khiển được tạo ra bởi phương trình (2.49) sẽ không gây ra bất cứ ảnh hưởng nào tới tất cả biến trạng thái của thiết bị. Hơn nữa phương trình (2.49) yêu cầu toàn bộ các biến trạng thái của hệ thống phải được đo lường, và có khả năng hồi tiếp ngược trở lại bộ điều khiển. Do đó bộ điều khiển bao gồm các cảm biến vật lý, thứ sẽ thực hiện nhiệm vụ đo các biến trạng thái, các thiết bị cơ khí và điện tử được gọi là bộ dẫn động, nó cung cấp ngõ vào cho hệ điều khiển dựa trên các ngõ ra mong muốn và định luật điều khiển (2.49). Bộ điều khiển hiện đại luôn luôn sử dụng mạch điện tử kỹ thuật số để thực hiện định luật điều khiển trong phần cứng. 
\paragraph{}Ma trận tăng cường bộ điều khiển, $K, K_{d}$ và $K_{n}$ là các tham số thiết kế của hệ thống điều khiển được mô tả bởi các phương trình (2.49) và (2.51). Cần phải chú ý rằng bậc của hệ thống vòng kín thông tin hồi tiếp toàn trạng thái là tương tự như với thiết bị. Biểu đồ của hệ thống điều khiển hồi tiếp toàn trạng thái được trình bày bên dưới:
\begin{figure}[H]
\begin{center}
	\includegraphics[scale=.35]{image/hinh16.png}
	\caption{Sơ đồ nguyên lý của hệ thống vòng kín thông tin hồi tiếp toàn trạng thái chung}
\end{center}
\end{figure}
\paragraph{}Ta xem xét hệ thống điều khiển có véc-tơ trạng thái mong muốn $x_{d}(t) = 0$, hệ thống này còn được gọi là bộ ổn định. Ngoài ra, để đơn giản hoá quá trình tính toán, chúng ta giả sử rằng tất cả các phép đo là hoàn hảo và không có bất kỳ lỗi nào xảy ra trong việc mô hình hoá thiết bị bởi hai phương trình (2.50) và (2.51). Như vậy, tất cả các yếu tố ngõ vào không mong muốn dưới dạng nhiễu đều bị loại bỏ, tức là $x_{n}(t) = 0$. Do đó phương trình của định luật điều khiển giảm xuống còn: 
\begin{equation}
u(t) = -Kx(t)
\end{equation}
\begin{figure}[H]
	\begin{center}
		\includegraphics[scale=.35]{image/hinh17.png}
		\caption{Sơ đồ nguyên lý của bộ ổn định hồi tiếp toàn trạng thái (hệ thống điều khiển có véc-tơ trạng thái mong muốn bằng 0) không có nhiễu}
	\end{center}
\end{figure}
\paragraph{}Thay phương trình(2.52) vào hai phương trình (2.50) và (2.51) ta có trạng thái vòng kín và phương trình ngõ ra của bộ ổn định như sau:
\begin{equation}
x^{(1)}(t) = (A-BK)x(t) 
\end{equation} 
\begin{equation}
y(t) = (C-DK)x(t) 
\end{equation}
\paragraph{}Phương trình (2.53) và (2.54) cho thấy bộ ổn định là hệ thống đồng nhất, được mô tả bởi ma trận hệ số trạng thái vòng kín A$_{CL}$ = A-BK, B$_{CL}$ = 0, C$_{CL}$ = C-DK và D$_{CL}$ = 0. Các cực vòng kín là giá trị riêng của ma trận A$_{CL}$. Dó đó bằng cách chọn ma trận tăng cường bộ ổn định K, chúng ta có thể đặt lại các cực vòng kín tại những vị trí mong muốn. Đối với thiết bị bậc \textit{n} với \textit{r} ngõ vào, kích thước của ma trận K là \textit{r}$\times$\textit{n}. Từ đó chúng ta sẽ nắm trong tay toàn bộ các tham số thiết kế vô hướng. Đối với hệ thống đa ngõ vào, số lượng của tham số thiết kế sẽ nhiều hơn số khả năng chọn lựa vị trí cực \textit{n}. 
\subsection{Thiết kế hệ thống ổn định gán cực cho ngõ vào đơn}
\paragraph{}Có một vấn đề gặp phải đó là, ngay cả với một thiết bị có ngõ vào đơn, bậc hai thì việc tính toán cho ma trận tăng cường bộ ổn định K bằng tay khá là rắc rối và có thể vượt quá khả năng khi mà bậc của thiết bị tăng vượt qua bậc ba. May mắn thay, nếu thiết bị ở \textit{dạng bộ điều khiển đôi}, thì việc tính toán đối với các thiết bị ngõ vào đơn lúc này sẽ được đơn giản hoá đi rất nhiều. Hãy xem xét một thiết bị ngõ vào đơn có bậc \textit{n}, \textit{dạng bộ điều khiển đôi} của nó sẽ như sau: 
\begin{equation}
A = 
\begin{bmatrix}
-a_{n-1} & -a_{n-2} & -a_{n-3} & ... & -a_{1} & -a_{0}\\
1 & 0 & 0 & ... & 0 & 0\\
0& 1 & 0 & ... & 0 & 0 \\
0 & 0 & 1 & ... & 0 & 0\\
. & . & . & ... & . & . \\
. & . & . & ... & . & . \\
0 & 0 & 0 & ... & 1 & 0 \\
0 & 0 & 0 & ... & 0 & 1
\end{bmatrix}
;
B = 
\begin{bmatrix}
1\\
0\\
0\\
0\\
.\\
.\\
0\\
0
\end{bmatrix}
\end{equation}
\paragraph{} Với $a_{0},..., a_{n-1}$ là các hệ số của đa thức đặc trưng của thiết bị $|sI - A| = s^{n} + a_{n-1}s^{n-1} + ... a_{1}s + a_{0}$. Ma trận tăng cường bộ ổn định hồi tiếp toàn trạng thái là một véc-tơ hàng gồm n tham số chưa biết được cho bởi:
\begin{equation}
K = 
\begin{bmatrix}
K1; K2; ...; K_{n}
\end{bmatrix}
\end{equation}
\paragraph{}Nó đòi hỏi phải đặt các cực vòng kín sao cho ma trận đặc đặc trưng vòng kín như sau:
\begin{equation}
|sI - A_{CL} = |sI - A + BK| = s^{n} + \alpha_{n-1}s^{n-1} + \alpha_{n-2}s^{n-2} ... + \alpha_{1}s + \alpha_{0}
\end{equation}
\paragraph{}Trong đó ma trận động lực học trạng thái vòng kín, $A_{CL} = A - BK$, được định nghĩa như sau:
\begin{equation}
A_{CL} = 
\begin{bmatrix}
(-a_{n-1} - K_{1}) & (-a_{n-2} - K_{2}) & ... & (-a_{1} - K_{n-1}) & (-a_{0} - K_{n})\\
1 & 0 & ... & 0 & 0\\
0& 1 & ... & 0 & 0\\
.&.&...&.&.\\
0&0&...&1&0\\
0&0&...&0&1
\end{bmatrix}
\end{equation}
\paragraph{}Lưu ý rằng hệ thống vòng kín cũng ở \textit{dạng bộ điều khiển đôi}. Do đó từ ma trận (2.58), các hệ số của đa thức đặc trưng vòng kín phải được diễn tả như sau:
\begin{equation}
\alpha_{n-1} = a_{n-1} + K_{1}; \alpha_{n-2} = a_{n-2} + K_{2}; ...; \alpha_{1} = a_{1}+K_{n-1}; \alpha_{0} = a_{0} + K_{n}
\end{equation} 
\paragraph{}Hoặc, các tham số ổn định chưa biết được tính đơn giản như sau:
\begin{equation}
K_{1} = \alpha_{n-1} - a_{n-1}; K_{2} = \alpha_{n-2} - a_{n-2}; ...; K_{n-1} = \alpha_{1} - a_{1}; K_{n} = \alpha_{0} - a_{0}
\end{equation}
\paragraph{}Ở dạng véc-tơ, phương trình (2.60) có thể được biểu diễn thành:
\begin{equation}
K = \alpha - a
\end{equation}
\paragraph{}Với $\alpha = [\alpha_{n-1}; \alpha_{n-2}; ...; \alpha_{1}; \alpha_{0}]$ và $a = [a_{n-1}; a_{n-2}; ...; a_{1}; a_{0}]$. Nếu biểu diễn không gian-trạng thái của thiết bị không ở dạng bộ điều khiển đôi, thì có thể sử dụng một phép biến đổi trạng thái để chuyển đổi thiết bị sang dạng bộ điều khiển đôi như sau:
\begin{equation}
x'(t) = Tx(t); A' = TAT^{-1}; B' = TB
\end{equation}
\paragraph{}Với $x'(t)$ là véc-tơ trạng thái của thiết bị ở dạng bộ điều khiển đôi, $x(t)$ là véc-tơ trạng thái ban đầu và $T$ là ma trận chuyển đổi trạng thái. Do đó định luật kiểm soát của bộ ổn định ngõ vào đơn (2.52) có thể được biểu diễn như sau: 
\begin{equation}
u(t) = -Kx(t) = -KT^{-1}x'(t)
\end{equation}
\paragraph{}Vì $KT^{-1}$ là ma trận tăng cường bộ ổn định khi thiết bị ở dạng bộ điều khiển đôi, nó phải được cho bởi phương trình (2.61):
\begin{equation}
KT^{-1} = \alpha - a
\end{equation}
\paragraph{}Hoặc:
\begin{equation}
K = (\alpha - a)T
\end{equation}
\paragraph{}Ma trận kiểm tra khả năng điều khiển của thiết bị ở không gian - trạng thái ban đầu được biểu diễn như sau: 
\begin{equation}
P = [B; AB; A^{2}B; ...; A^{n-1}B]  
\end{equation}
\paragraph{}Sử dụng phép biến đổi nghịch đảo, thay $B = T^{-1}B'$ và $A = T^{-1}A'T$ vào phương trình (2.66):
\begin{align} 
P 
&= [T^{-1}B'; (T^{-1}A'T)T^{-1}B'; (T^{-1}A'T)^{2}T^{-1}B'; ...; (T^{-1}A'T)^{n-1}T^{-1}B']\\
\nonumber
&= T^{-1}[B'; A'B'; (A')^{2}B'; ...; (A')^{n-1}B'] = T^{-1}P'
\end{align}
\paragraph{}Với $P'$ là ma trận kiểm tra khả năng điều khiển của thiết bị ở \textit{dạng bộ điều khiển đôi}. Nhân cả 2 vế của phương trình (2.67) với ma trận chuyển đổi trạng thái $T$, sau đó tiếp tục nhân cả 2 vế phương trình thu được với $P^{-1}$, chúng ta nhận được biểu thức sau cho $T$: 
\begin{equation}
T = P'P^{-1}
\end{equation}
\paragraph{}Có thể dễ dàng hiểu rằng $P'$ là ma trận tam giác trên (gọi như vậy vì tất cả các phần tử bên dưới đường chéo chính của nó đều bằng 0):
\begin{equation}
P' = 
\begin{bmatrix}
1&-a_{n-1}&-a_{n-2}&...&-a_{2}&-a_{1}\\
0&1&-a_{n-1}&...&-a_{3}&-a_{2}\\
0&0&1&...&-a_{4}&-a_{3}\\
.&.&.&...&.&.\\
.&.&.&...&.&.\\
0&0&0&...&1&-a_{1}\\
0&0&0&...&0&1
\end{bmatrix}
\end{equation}
\paragraph{}Thay phương trình (2.68) vào (2.65), ma trận tăng cường bộ ổn định do đó trở thành: 
\begin{equation}
K = (\alpha - a)P'P^{-1}
\end{equation}
\paragraph{}Phương trình (2.70) được gọi là \textit{công thức không gian đặt cực Ackermann}. Đối với thiết bị ngõ vào đơn được xem xét ở đây, $P$ và $P'$ đều là ma trận vuông kích thước (\textit{n} $\times$ \textit{n}). 
\paragraph{}Áp dụng vào thiết kế hệ thống điều chỉnh thông tin phản hồi trạng thái đầy đủ cho mô hình xe hai bánh tự cân bằng. Từ phương trình biểu diễn không gian - trạng thái tuyến tính của hệ thống, ma trận hệ số trạng thái được biểu diễn như sau: 
\begin{equation}
\begin{bmatrix}
\dot{x_{W}}\\
\ddot{x_{W}}\\
\dot{\theta}\\
\ddot{\theta}
\end{bmatrix}
=
\begin{bmatrix}
0&1&0&0\\
0&0&(g-\dfrac{4}{3}L\dfrac{\ddot{M_{P}}g}{X}) &0\\
0&0&0&1\\
0&0&\dfrac{M_{P}g}{X}&0
\end{bmatrix}
\begin{bmatrix}
x_{W}\\
\dot{x_{W}}\\
\theta\\
\dot{\theta}
\end{bmatrix}
+
\begin{bmatrix}
0\\
\dfrac{4LY}{3X} - \dfrac{1}{M_{P}L}\\
0\\
\dfrac{-Y}{X}
\end{bmatrix}
C_{\theta}
\end{equation}
\begin{equation}
\begin{bmatrix}
x_{W}\\
\theta
\end{bmatrix}
= 
\begin{bmatrix}
1&0&0&0\\
0&0&1&0
\end{bmatrix}
\begin{bmatrix}
x_{W}\\
\dot{x_{W}}\\
\theta\\
\dot{\theta}
\end{bmatrix}
+
\begin{bmatrix}
0\\
0
\end{bmatrix}
C_{\theta}
\end{equation}
\paragraph{}Với:
\begin{equation}
X = \left[\dfrac{4}{3}M_{P}L-\dfrac{M_{P}}{2M_{W}+M_{P}}(M_{P}L+M_{W}R)\right]
\end{equation}
\begin{equation}
Y = \dfrac{M_{P}}{2M_{W}+M_{P}}\dfrac{1}{R}+\dfrac{1}{L}
\end{equation}
\paragraph{}Ngõ vào đơn, $u(t)$, là một lực tác dụng lên thiết bị theo phương nằm ngang và hai ngõ ra là vị trí góc của nghiêng của thiết bị, $\theta(t)$, và vị trí ngang của sàn xe, $x(t)$. Bốn véc-tơ trạng thái ứng với các bậc tự do của thiết bị lần lượt là $x(t) = [\theta(t); x(t); \theta^{1}(t); x^{1}(t)]^{T}$. Áp dụng các tham số thực tế: $M_{W} =$ 7 kg, $M_{P} =$ 60 (kg), $R$ = 0.2 (m), $L$ = 1 (m) và $g$ = 9.8 (m/s$^{2}$). Ma trận A và B lúc này trở thành:
\begin{equation}
A = 
\begin{bmatrix}
0&1&0&0\\
0&0&-15.962&0\\
0&0&0&1\\
0&0&19.329&0
\end{bmatrix};
B = 
\begin{bmatrix}
0\\
0.2037\\
0\\
-0.1653
\end{bmatrix}
\end{equation}
\paragraph{}Nhập giá trị của ma trận A và B vào MATLAB bằng cú pháp:
\begin{lstlisting}
>> A =  [0 1 0 0; 0 0 -15.9620 0;0 0 0 1;0 0 19.3290 0]
<enter>
>> B = [0;0.2037;0;-0.1653] 
<enter>
\end{lstlisting}
\paragraph{}Sử dụng ma trận kiểm tra khả năng điều khiển để kiểm tra xem thiết bị mà chúng ta xây dựng có thể điều khiển được hay không bằng cách dùng lệnh $ctrb$ có sẵn trong MATLAB:
\begin{lstlisting}
>> P = ctrb(A,B) 
<enter>
\end{lstlisting}
\paragraph{}Ta thu được ma trận P:
\begin{equation}
P =
\begin{bmatrix}
0&0.2037&0&2.6385\\
0.2037&0&2.6385&0\\
0&-0.1653&0&-3.1951\\
-0.1653&0&-3.1951&0
\end{bmatrix}
\end{equation}
\paragraph{}Định thức của ma trận kiểm tra khả năng điều khiển P bằng:
\begin{lstlisting}
>> det(P)
<enter>
ans = 
      0.0461
\end{lstlisting}
\paragraph{}Vì |P| = 0.0461 $\neq$ 0, như vậy thông qua việc kiểm tra định thức của ma trận kiểm tra khả năng điều khiển đã cho ta biết rằng thiết bị với các thông số đã cho bên trên là hoàn toàn có thể điều khiển được. 
\section{Các phương pháp xử lý tín hiệu từ cảm biến}
\paragraph{}Như đã trình bày, giá trị ngõ ra được quan tâm hàng đầu của thiết bị hai bánh tự cân bằng chính là góc giữa tay lái thẳng đứng hoặc của sàn xe hợp với chiều trọng lực. Hiện nay có nhiều loại cảm biến có thể được dùng để đo góc như encoder, inclinometer, resolver,..., cùng nhiều phương pháp như đo trực tiếp góc dựa trên khoảng cách, đo gián tiếp không tiếp xúc bằng hiệu ứng áp điện,... 
\begin{center}
	\begin{figure}[H]
		\begin{center}
			\includegraphics[scale=.3]{image/hinh18.png}
			\caption{Một số cảm biến đo góc phổ biến hiện nay}
		\end{center}
	\end{figure}
\end{center}
Trong đề tài này, thiết bị sử dụng hai loại cảm biến để xác định góc giữa sàn xe với trọng lực theo phương pháp không tiếp xúc là con quay hồi chuyển để đo gia tốc góc và gia tốc kế để đo góc nghiêng tĩnh. Tín hiệu thu được từ các cảm biến khác nhau mang những ý nghĩa khác nhau, việc kết hợp các tín hiệu này để nhận được tín hiệu cần thiết duy nhất là vô cùng quan trọng, đồng thời trong quá trình truyền dẫn và xử lý việc xảy ra nhiễu gây ảnh hưởng đến số liệu là điều hoàn toàn không thể tránh khỏi. Để làm mịn và giảm thiểu tối đa sai lệch do nhiễu gây ra cũng như để kết hợp tín hiệu từ hai cảm biến con quay hồi chuyển và gia tốc kế, người ta thường sử dụng các bộ lọc được xây dựng từ phần cứng hoặc phần mềm, ví dụ như lọc bộ phụ thông tần, lọc thích nghi, lọc trung bình,... 
\subsection{IMU}
\paragraph{}Đầu tiên, ta nói đến hai khái niệm IMU và DOF. IMU, viết tắt của Inertial Measurement Unit - Hệ thống đo lường quán tính, là một con chip có nhiệm vụ đo những chuyển động nghiêng, xoay, lắc,... Một module IMU thường gồm có 2 loại cảm biến: cảm biến gia tốc (gia tốc kế) và cảm biến quay (con quay hồi chuyển). 
\begin{itemize}
	\item Cảm biến gia tốc (Accelerometer): là một cảm biến đo lường "gia tốc" của bản thân module, nhưng thực ra là đo lực trên mỗi đơn vị khối lượng ($F = ma$, vậy nên $a = F/m$). Cảm biến gia tốc thường sẽ có 3 trục XYZ ứng với 3 chiều không gian (loại 1 và 2 trục ít được sử dụng), được sử dụng để đo trọng lực hoặc góc nghiêng. Đối với nhiệm vụ đo góc nghiêng, sẽ được xác định thông qua trục X, trong khi đó trục Y và trục Z lại kém nhạy cảm hơn với những thay đổi nhỏ của góc so với trục X cũng như chúng không bị phụ thuộc vào hướng nghiêng. 
	\begin{center}
		\begin{figure}[H]
			\begin{center}
				\includegraphics[scale=1.5]{image/hinh19.png}
				\caption{Có thể nhận thấy được chiều của trọng lực thông qua trục X}
			\end{center}
		\end{figure}
	\end{center}
	\item Cảm biến quay (Gyroscope): là một loại cảm biến đo tốc độ quay của nó quanh một trục, hay còn gọi là tốc độ góc. Kết quả đọc được từ cảm biến quay sẽ bằng "0" khi nó đứng im và có giá trị âm hoặc dương tuỳ theo chiều quay của nó. Tương tự như với cảm biến gia tốc, cảm biến quay cũng thường có 3 trục XYZ. 
		\begin{center}
		\begin{figure}[H]
			\begin{center}
				\includegraphics[scale=1.7]{image/hinh20.png}
				\caption{Giá trị đọc được từ cảm biến quay tuỳ theo chiều quay}
			\end{center}
		\end{figure}
	\end{center}
\end{itemize}
\paragraph{}Một module IMU đầy đủ sẽ được gọi là 6-DOF (Degrees Of Freedom), tức là 6 trục độc lập (3 của cảm biến gia tốc và 3 của cảm biến quay). Ngoài ra nhằm mục đích phục vụ những dự án phức tạp hơn như là để điều hướng máy bay hoặc robot, người ta cần phải sử dụng đến các IMU 9-DOF (thêm một cảm biến từ trường 3 trục - Magnetometer - hoạt động gần giống như một la bàn để định hướng), hoặc 10-DOF (thêm một áp kế - Barometer - dùng để đo độ cao), thậm chí là 11-DOF (thêm module GPS để xác định vị trí theo thời gian thực). Tuy nhiên phần lớn các ứng dụng hiện nay chỉ cần một IMU 6-DOF vẫn có thể đáp ứng hoàn hảo. 
\paragraph{}Một số ứng dụng của chip IMU trong đời sống hàng ngày:
\begin{itemize}
	\item Xe hai bánh tự cân bằng
	\item Máy bay không người lái
	\item Tay cầm chơi game
	\item Hệ thống chống rung cho máy ảnh, máy quay phim
	\item Đồng hồ thông minh
	\item vv...vv
\end{itemize}
\paragraph{}IMU - cũng giống như nhiều cảm biến khác - cần được hiểu chỉnh trước khi sử dụng nếu muốn có một kết quả đáng tin cậy. Một số vấn đề thường gặp như:
\begin{itemize}
	\item Đối với cảm biến gia tốc: luôn có offset trên mỗi trục làm cho các giá trị đo thường lệch đi so với thực tế một chút. Ngoài ra, giá trị đo được theo gia tốc kế thường khá nhiễu do nó rất nhạy với các rung động cơ khí nhỏ khiến cho giá trị tức thời của nó không đáng tin cậy và khó đọc. 
	\item Đối với cảm biến quay: cũng như cảm biến gia tốc, cảm biến quay cũng có offset (hay còn gọi là bias) làm lệch các giá trị đo. Một vấn đề khác nữa có thể gặp phải của cảm biến quay là driff (độ trôi tín hiệu), có nghĩa là bias cũng thay đổi chậm theo thời gian. Nguyên nhân là do các tác động cơ khí, rung động tác động lên cảm biến, sau một thời gian các giá trị này tích luỹ lên đáng kể, làm cho giá trị đo góc không còn được chính xác. Dù vậy, điểm mạnh là cảm biến quay ít nhiễu hơn cảm biến gia tốc, tức là giá trị tức thời của nó đáng tin cậy hơn. 
\end{itemize}
\subsection{Đọc dữ liệu từ cảm biến}
\paragraph{}Bước đầu tiên là đọc ngõ vào tương tự (analog) thông qua bộ chuyển đổi tín hiệu tương tự sang tín hiệu số (ADC) cho từng cảm biến và lọc ra các dữ liệu hữu ích cho hệ thống. Việc này yêu cầu có sự điều chỉnh đối với offset và scale:
\begin{itemize}
	\item Rất dễ dàng để tìm ra được offset đối với mỗi loại cảm biến: quan sát giá trị nguyên mà các cảm biến đọc được khi nó đang ở vị trí cân bằng/hoặc đứng yên. Nếu giá trị này biến đổi liên tục xung quanh một giá trị nguyên nào đó thì lấy giá trị trung bình của nó. Offset phải là một dạng biến integer có dấu (hoặc hằng số). 
	\item Khác với offset, scale tuỳ thuộc vào từng loại cảm biến. Nó là một hệ số được nhân thêm vào để có thể đạt được giá trị mong muốn. Hệ số này thường được tìm thấy trong bảng dữ liệu do nhà sản xuất cung cấp hoặc thông qua thực nghiệm. Đôi khi nó còn được gọi là hằng số cảm biến, độ tăng hoặc độ nhạy. Scale phải là biến dạng float (hoặc hằng số).
\end{itemize}
\subsection{Phép đo lý tưởng}
\paragraph{}Để điều khiển được thiết bị, sẽ rất tốt khi biết được cả hai thông tin là góc và tốc độ góc của nó. Những thông tin này có thể là cơ sở cho thuật toán điều khiển góc PD (tỷ lệ/đạo hàm), đây là một loại thuật toán đã được chứng minh là hoạt động tốt cho loại hệ thống cân bằng chúng ta đang xây dựng. Có thể hình dung như sau:
\begin{equation}
Motor\_Output = K_{p} \times (Angle) + K_{d} \times (Angular\_Velocity) 
\end{equation}
\paragraph{}Ý tưởng chung là thiết lập điều khiển này có thể được điều chỉnh bằng hai biến $K_{p}$ và $K_{d}$ để mang lại sự ổn định và hiệu suất đúng như ta mong muốn. Nó ít có khả năng bị vượt quá điểm cân bằng hơn so với bộ điều khiển chỉ theo tỷ lệ (Ví dụ: Nếu góc của thiết bị so với phương ngang là giá trị dương nhưng vận tốc góc lại mang giá trị âm, tức là nó đang quay trở lại vị trí cân bằng, động cơ quay chậm lại về phía trước).  
\newline
\begin{wrapfigure}{l}[0pt]{0.5\linewidth}
	\begin{center}
		\includegraphics[scale=.5]{image/hinh21.png}
		\caption{Sơ đồ điều khiển PD}
	\end{center}
\end{wrapfigure}
\newline
\newline
\paragraph{} Trên thực tế, sơ đồ điều khiển PD giống như việc thêm một lò xo và van điều tiết có thể điều chỉnh được tới thiết bị lái tự cân bằng. 
\newline
\newline
\newline
\newline
\newline
\newline
\subsection{Kết hợp dữ liệu từ các cảm biến (Ánh xạ cảm biến)}
\paragraph{}Như đã nói, việc kết hợp các dữ liệu (ở đây là tín hiệu dạng số) từ các cảm biến là một công việc hết sức quan trọng và cần phải được xử lý cẩn thận để có thể thu được các thông tin cần thiết và tin cậy phục vụ quá trình điều khiển thiết bị. 
	\begin{figure}[H]
		\begin{center}
			\includegraphics[scale=1.5]{image/hinh22.png}
			\caption{Bài toán kết hợp dữ liệu từ các cảm biến}
		\end{center}
	\end{figure}
\paragraph{}Có rất nhiều phương pháp để kết hợp dữ liệu và lọc các thông tin nhiễu. Tuy nhiên mỗi phương pháp lại có những ưu điểm và nhược điểm riêng biệt, tuỳ theo mục đích và nhu cầu chúng ta có thể lựa chọn phương pháp thích hợp và tối ưu nhất. Đa số các phương pháp hiện nay đều dựa trên lý thuyết bộ lọc kỹ thuật số, giúp chúng ta có thể lập trình, tuỳ chỉnh một cách dễ dàng mà không cần lo lắng về vấn đề thay đổi cũng như chi phí dành cho phần cứng. Một số định nghĩa thường được nhắc tới trong bộ lọc kỹ thuật số:
\begin{itemize}
	\item \textbf{Tích phân hoá} Có thể được hiểu đơn giản thông qua ví dụ sau: Có một chiếc ô tô đang di chuyển với một tốc độ cho trước, và chương trình của chúng ta là một chiếc đồng hồ mà cứ mỗi vài mili giây nó sẽ nhảy một nhịp. Để có được vị trí của chiếc ô tô sau mỗi nhịp đồng hồ, chúng ta sẽ lấy vị trí cũ và cộng thêm sự thay đổi về vị trí. Sự thay đổi về vị trí chính là tốc độ của ô tô nhân với thời gian đã trôi qua kể từ nhịp cuối cùng - chúng ta có thể biết giá trị này thông qua bộ đếm thời gian tích hợp sẵn trong vi điều khiển. Ứng dụng trong chương trình về hệ thống lái cân bằng: $angle$ $+$$= gyro * dt$.
	\item \textbf{Bộ lọc thông thấp} Mục đích của bộ lọc thông thấp là chỉ cho phép những tín hiệu thay đổi mang tính dài hạn được phép đi qua, còn các tín hiệu biến động ngắn hạn sẽ bị loại bỏ. Để thực hiện việc này bằng chương trình phần mềm, có một cách là buộc các tín hiệu thay đổi phải tăng lên một chút một trong những lần tiếp theo thông qua vòng lặp chương trình: $angle = 0.98 * angle + 0.02 * x\_acc$.
	\item \textbf{Bộ lọc thông cao} Cho phép các tín hiệu thay đổi mang tính ngắn hạn đi qua trong khi các tín hiệu dài hạn ổn định theo thời gian sẽ bị loại bỏ. Điều này có thể được sử dụng để giảm thiểu hoặc triệt tiêu hoàn toàn các tín hiệu trôi. 
	\item \textbf{Thời gian lấy mẫu} Khoảng thời gian trôi qua giữ mỗi vòng lặp chương trình. Tốc độ lấy mẫu thường được sử dụng là 100Hz, tương ứng với thời gian lấy mẫu là 0,01 giây. 
	\item \textbf{Hằng số thời gian} Hằng số thời gian của bộ lọc là khoảng thời gian tương đối của tín hiệu mà nó sẽ hoạt động. Đối với bộ lọc thông thấp, các tín hiệu dài hơn nhiều so với hằng số thời gian sẽ được đi qua một cách nguyên vẹn, trong khi đó các tín hiệu ngắn hơn sẽ bị loại bỏ. Điều ngược lại cũng đúng với bộ lọc thông cao. Hằng số thời gian, $\uptau$, của bộ lọc thông thấp: $y = \alpha * y + (1-\alpha) * x$, chạy trong một vòng lặp chương trình với thời gian lấy mẫu, $dt$, có thể dễ dàng xác định như sau: 
	\begin{equation}
	\uptau = \dfrac{\alpha*dt}{1-\alpha} \leftrightarrow \alpha = \dfrac{\uptau}{\uptau + dt}
	\end{equation}
	Như vậy chúng ta có thể chọn được hệ số $\alpha$ thích hợp cho bộ lọc.
\end{itemize}
\paragraph{}\textbf{1. Phương pháp trực tiếp}
\begin{figure}[H]
	\begin{center}
		\includegraphics[scale=1.5]{image/hinh23.png}
		\caption{Phương pháp trực tiếp}
	\end{center}
\end{figure}
\paragraph{}Ưu điểm:
\begin{itemize}
	\item Trực tiếp
	\item Dễ dàng để lập trình
	\item Cảm biến quay cho dữ liệu tốc độ góc nhanh và chính xác
\end{itemize}
\paragraph{}Nhược điểm:
\begin{itemize}
	\item Dữ liệu bị nhiễu
	\item Trục X của gia tốc kế sẽ đọc bất kỳ gia tốc theo phương ngang nào như là một sự thay đổi của góc (Hãy tưởng tượng sàn thiết bị được đặt nằm ngang, nhưng các động cơ đang tăng tốc làm cho nó tiến về phía trước, gia tốc kế không thể phân biệt được điều này với lực hấp dẫn).
\end{itemize}
\paragraph{}\textbf{2. Phương pháp lọc nhanh}
\begin{figure}[H]
	\begin{center}
		\includegraphics[scale=1.5]{image/hinh24.png}
		\caption{Phương pháp lọc nhanh}
	\end{center}
\end{figure}
\paragraph{}Sử dụng một bộ lọc thông thấp để lọc dữ liệu thu được từ cảm biến gia tốc. Có thể thực hiện đơn giản như việc lấy trung bình các mẫu: 
\begin{equation}
angle = (0.75) * (angle) + (0.25) * (x\_acc)
\end{equation}
\paragraph{}Ưu điểm:
\begin{itemize}
	\item Trực tiếp
	\item Dễ dàng để lập trình
	\item Lọc bỏ các tín hiệu gia tốc ngắn hạn theo phương ngang. Chỉ các tín hiệu gia tốc dài hạn (trọng lực) mới được đi qua bộ lọc
\end{itemize}
\paragraph{}Nhược điểm:
\begin{itemize}
	\item Việc đo góc sẽ bị trễ dần theo thời gian. Càng lọc nhiều thì sẽ càng bị trễ. Độ trễ chính là một tác nhân trực tiếp tới sự ổn định của thiết bị
\end{itemize}
\paragraph{}\textbf{3. Phương pháp đơn cảm biến}
\begin{figure}[H]
	\begin{center}
		\includegraphics[scale=1.5]{image/hinh25.png}
		\caption{Phương pháp đơn cảm biến}
	\end{center}
\end{figure}
\paragraph{}Thay vì phải sử dụng hai thông tin góc và tốc độ góc thu được trực tiếp từ hai cảm biến, chúng ta chỉ cần dùng duy nhất một cảm biến là con quay để có được cả hai thông tin này bằng kiến thức vật lý hết sức đơn giản: \textit{khoảng cách} = \textit{tốc độ góc} $\times$ \textit{thời gian}. Có thể được lập trình dễ dàng bằng đoạn mã như sau (yêu cầu phải biết được thời gian giữa các lần cập nhật biến, \textit{dt}):
\begin{equation}
angle = angle + gyro * dt
\end{equation}
\paragraph{}Ưu điểm:
\begin{itemize}
	\item Chỉ cần đọc dữ liệu từ một cảm biến
	\item Tốc độ nhanh, độ trễ không còn là vấn đề
	\item Không bị lệ thuộc vào gia tốc theo phương ngang
	\item Dễ dàng để lập trình
\end{itemize}
\paragraph{}Nhược điểm:
\begin{itemize}
	\item Độ trôi của con quay hồi chuyển là điều rất đáng lo ngại. Nếu như nó không thể đọc được chính xác giá trị bằng 0 khi đang ở trạng thái đứng yên (và chắc chắn là nó sẽ không thể đọc được chính xác), một phần nhỏ giá trị góc sai lệch sẽ được thêm dần vào sau mỗi lần cập nhật cho đến khi nó hoàn toàn bị khác xa so với góc thực tế 
\end{itemize}
\paragraph{}\textbf{4. Phương pháp lọc thích nghi - Bộ lọc Kalman}
\begin{figure}[H]
	\begin{center}
		\includegraphics[scale=1.5]{image/hinh26.png}
		\caption{Phương pháp lọc thích nghi}
	\end{center}
\end{figure}
\paragraph{}Lý thuyết lọc Kalman được từ những năm 1960 bởi giáo sư R. E. Kalman để thu thập và kết hợp linh động các thông tin từ các cảm biến thành phần. Một cách khái quát, bộ lọc Kalman là một tập hợp các phương trình truy hồi cho phép ta ước lượng trạng thái của một quá trình theo tiêu chuẩn bình phương nhỏ nhất. Nó rất hiệu quả trong việc ước lượng các trạng thái trong quá khứ, hiện tại và tương lai ngay cả khi chưa biết chính xác mô hình của hệ thống. 
\paragraph{}Đối với bộ lọc Kalman, thuật ngữ "lọc" không được hiểu theo nghĩa giống như các bộ lọc thông thường. Bộ lọc Kalman là một giải thuật xử lý dữ liệu hồi quy tối ưu dựa vào việc hợp nhất tất cả các thông tin được cung cấp ở ngõ vào để đưa ra thông tin ngõ ra tin cậy nhất. Bộ lọc có khả năng loại bỏ các nhiễu trắng, nhiều ồn, nhiễu đo mà nó nhận được ở ngõ vào dựa theo các thống kê trước đó và hiệu chỉnh lại bằng các giá trị đo hiện tại. 
\paragraph{}Trong thiết bị lái tự cân bằng, tín hiệu cảm biến đi vào bộ lọc gồm hai tín hiệu: góc từ cảm biến gia tốc kế và tốc độ góc từ cảm biến con quay. Hai tín hiệu này sẽ hỗ trợ và xử lý lẫn nhau thông qua mô hình vật lý: \textit{tốc độ góc} = \textit{đạo hàm/vi phân của góc}. Tín hiệu ngõ ra của bộ lọc là tín hiệu đã được xử lý và loại bỏ nhiễu. Nhờ có cơ chế cập nhật sai lệch bias tại mỗi thời điểm tính toán nên giá trị góc nghiêng ước lượng được ổn định và chính xác. 
\paragraph{}Ví dụ mô phỏng về bộ lọc Kalman: Gọi $u_{0} = 100.0$ là giá trị thực của tín hiệu, cũng là giá trị mà ta mong muốn thu được vì $u_{0}$ là hằng số nếu như không có bị ảnh hưởng bởi nhiễu. Sử dụng một biến $e$ để mô phỏng giá trị nhiễu, giá trị này được tạo ra bằng hàm random với biên độ bằng $\pm100\%$ $u_{0}$. Lúc này tín hiệu thu được sẽ là $u = u_{0} + e$.
\begin{figure}[H]
	\begin{center}
		\includegraphics[scale=0.2]{image/hinh27.png}
		\caption{Tín hiệu thu được khi chưa cho xử lý qua bộ lọc Kalman}
	\end{center}
\end{figure}
\paragraph{}Như đã nói trước ở phần trên, trong thực tế $u_{0}$ là giá trị chúng ta không biết, việc sử dụng bộ lọc sẽ phải giúp ta loại bỏ được các thành phần nhiễu, khi đó giá trị đo được phải gần đường $u_{0} = 100$ hơn. Vì đây là ví dụ để mô phỏng nên giá trị $u_{0}$ cần được biết trước để có thể kiểm chứng tính đúng đắn của kết quả thu được trước và sau khi được lọc (bằng cách trộn $u_{0}$ với tín hiệu nhiễu $e$). Giá trị thu được sau khi đi qua bộ lọc Kalman sẽ là $u\_kalman$. 
\begin{figure}[H]
	\begin{center}
		\includegraphics[scale=0.2]{image/hinh28.png}
		\caption{Tín hiệu thu được khi qua một tầng lọc Kalman}
	\end{center}
\end{figure}
\begin{figure}[H]
	\begin{center}
		\includegraphics[scale=0.2]{image/hinh30.png}
		\caption{Tín hiệu thu được khi qua bốn tầng lọc Kalman}
	\end{center}
\end{figure}
\paragraph{}Nhận xét: Có những lúc tín hiệu nhiều dồn ra biên cực đại ($\pm100$ $u_{0}$), nhưng giá trị $u\_kalman$ vẫn bám khá sát với đường $u_{0}$. Như vậy có thể thấy thuật toán lọc Kalman tỏ ra khá hiệu quả. Ngoài ra chúng ta còn có thể ghép tầng các bộ lọc để có thể thu được kết quả chính xác hơn, tuy nhiên sẽ bị đáp ứng trễ một tầng. 

\paragraph{}Ưu điểm:
\begin{itemize}
	\item Được cho là bộ lọc lý tưởng về mặt lý thuyết để kết hợp các tín hiệu nhiễu của cảm biến, đưa ra kết quả chính xác, đáng tin cậy
	\item Tham gia vào việc tính toán các thuộc tính vật lý đã biết của hệ thống
\end{itemize} 
\paragraph{}Nhược điểm:
\begin{itemize}
	\item Phức tạp về mặt toán học, đòi hỏi phải có một số kiến thức về đại số tuyến tính
	\item Có rất nhiều dạng khác nhau cho các trường hợp khác nhau
	\item Khó khăn trong việc lập trình, phải phụ thuộc vào các thư viện có sẵn
	\item Tốn rất nhiều thời gian xử lý
\end{itemize}
\paragraph{}\textbf{5. Phương pháp lọc bổ phụ thông tần}
\begin{figure}[H]
	\begin{center}
		\includegraphics[scale=1.5]{image/hinh29.png}
		\caption{Phương pháp lọc bổ phụ thông tần - Bộ lọc Complementary}
	\end{center}
\end{figure}
\paragraph{}Ý tưởng chính của bộ lọc Complementary là lấy tín hiệu chuyển động chậm từ gia tốc kế và tín hiệu chuyển động nhanh từ con quay hồi chuyển và kết hợp chúng lại với nhau. Gia tốc kế sẽ đưa ra các giá trị định hướng tốt trong điều kiện tĩnh, trong khi đó con quay hồi chuyển lại cung cấp các giá trị tốt về độ nghiêng trong điều kiện động. Vì vậy, bằng việc chuyển tín hiệu gia tốc kế qua bộ lọc thông thấp và tín hiệu con quay hồi chuyển qua bộ lọc thông cao và kết hợp chúng lại với nhau có thể đưa ra tỷ lệ cuối cùng. Điểm mấu chốt ở đây chính là tần số đáp ứng của bộ lọc thông thấp và thông cao cộng lại bằng 1 tại tất cả dải tần, nghĩa là tại bất kỳ thời điểm xác định nào đó tín hiệu hoàn chỉnh cuối cùng có thể là tín hiệu hạ tần hoặc cao tần. 
\paragraph{}Phương trình cho bộ lọc thông thấp sử dụng cho các góc thu được từ gia tốc kế:
\begin{equation}
y_{n} = (1 - \alpha)x_{n} + y_{n-1}\alpha 
\end{equation}
\paragraph{}Phương trình cho bộ lọc thông cao sử dụng cho các góc thu được từ con quay hồi chuyển:
\begin{equation}
y_{n} = (1-\alpha)y_{n-1} + (1-\alpha)(x_{n} - x_{n-1})
\end{equation}
\paragraph{}Trong đó:
\begin{itemize}
	\item  $x_{n}$ là giá trị yaw/pitch/roll nhận được từ gia tốc kế/con quay hồi chuyển
	\item $y_{n}$ là giá trị yaw/pitch/roll cuối cùng đã được lọc và được sử dụng để phục vụ các bước tiếp theo trong chương trình
	\item $n$ là chỉ số mẫu hiện tại
	\item $\alpha$ liên quan đến hằng số thời gian. Nó xác định một cột mốc mà ở đó việc đọc dữ liệu từ gia tốc kế sẽ bị dừng lại để chuyển sang đọc dữ liệu từ con quay hồi chuyển và ngược lại. Nó kiểm soát số lượng giá trị hiện tại hoặc giá trị mới được thêm mà ngõ ra sẽ phụ thuộc vào tuỳ theo yêu cầu của chúng ta. Giá trị của $\alpha$ trong hai bộ lọc phải hoàn toàn giống nhau, và thường lớn hơn 0.5
\end{itemize}
\paragraph{}Như đã nói ở phần trên, hệ số $\alpha$ có thể được tính thông qua hằng số thời gian, $\uptau$, bằng công thức:
\begin{equation*}
	\alpha = \dfrac{\uptau}{\uptau +dt}
\end{equation*}
\paragraph{}Thay vì phải sử dụng hai bộ lọc thông thấp và thông cao riêng biệt, chúng ta có thể sử dụng một phương trình duy nhất sau để kết hợp và lọc tín hiệu từ hai cảm biến gia tốc kế và con quay hồi chuyển, phần nào giảm thiểu khối lượng công việc cũng như thời gian xử lý cho chương trình:
\begin{equation}
angle = (1-\alpha)(angle + gyro * dt) + (\alpha)(x\_acc)
\end{equation}
\paragraph{}Giá trị đọc được đầu tiên thu được từ quá trình tích phân hoá tín hiệu con quay hồi chuyển, giá trị đọc được tiếp theo là từ gia tốc kế. Trong trường hợp biến $gyro = 0$, góc sẽ hội tụ về góc xác định bởi gia tốc kế. Thử với một ví dụ sau:
\begin{equation}
angle = (0.98)*(angle + gyro*dt) + (0.02)*(x\_acc)
\end{equation}
\paragraph{}Ta có thể thấy bước $(angle +gyro *dt)$ chính là bước tích phân hoá tín hiệu con quay hồi chuyển. Sau đó nó được nhân với $(1 - \alpha) = (0.98)$, đây giống như một bộ lọc thông cao dựa trên góc ước tính của con quay hồi chuyển sau khi đã được lấy tích phân. Trong khi đó $(0.02)*(x\_acc)$ là bộ lọc thông thấp hoạt động với gia tốc kế. Hằng số thời gian ở cả hai bộ lọc là hoàn toàn như nhau. 
\paragraph{}Nếu bộ lọc đang chạy trong một vòng lặp 100 lần/giây, thì hằng số thời gian cho cả bộ lọc thông thấp và bộ lọc thông cao là:
\begin{equation}
\uptau = \dfrac{a * dt}{1-\alpha} = \dfrac{0.98 * 0.01}{0.02} = 0.49 
\end{equation}
\paragraph{}Phép tính trên giúp chúng ta xác định được đâu là biên thời gian mà tại đó chương trình điều khiển sẽ "tin tưởng" giá trị của cảm biến nào hơn. Ví dụ như với hằng số thời gian bằng 0.5 giây. Đối với các khoảng thời gian ngắn hơn nửa giây, kết quả tích phân hoá của con quay hồi chuyển sẽ được ưu tiên và các nhiễu trục ngang của gia tốc kế sẽ được loại bỏ. Tiếp theo trong khoảng thời gian dài hơn nửa giây, giá trị trung bình của gia tốc kế có hệ số "tin tưởng" lớn hơn so với con quay hồi chuyển, thứ mà khả năng cao đang xảy ra sự trôi tín hiệu tại thời điểm đó. 
\paragraph{}Việc thiết kế bộ lọc bổ phụ thông tần thường bắt đầu từ việc lựa chọn hằng số thời gian để từ đó sử dụng nó cho việc tính toán các hệ số còn lại của bộ lọc. Bước chọn hằng số thời gian chính là bước quan trọng giúp chúng ta có thể điều chỉnh được sự phản hồi của hệ thống. Nếu con quay hồi chuyển bị trôi trung bình 2 độ mỗi giây, chúng ta sẽ cần có một hằng số thời gian nhỏ hơn 1 giây để đảm bảo rằng không bao giờ bị lệch quá một vài độ theo cả hai hướng. Nhưng hằng số thời gian càng thấp thì sẽ có càng nhiều các tín hiệu nhiễu trục ngang của gia tốc kế được phép đi qua bộ lọc. Để chọn được các hệ số phù hợp thì tỷ lệ mẫu là một yếu tố vô cùng quan trọng. Cách duy nhất để có thể có được một hằng số thời gian tốt, đáp ứng được yêu cầu, giúp chúng ta thực sự điều chỉnh được hệ thống như ý muốn là phải thông qua quá trình thực nghiệm. 
\begin{figure}[H]
			\begin{center}
		\includegraphics[scale=0.2]{image/hinh31.png}

		\caption{Kết quả thử nghiệm với tỷ lệ mẫu 79 Hz, hằng số thời gian 0.62 giây và hệ số bộ lọc 0.98}
	\end{center}
\end{figure}
\paragraph{}Từ đồ thị bên trên, ta có thể thấy được cách mà bộ lọc bổ phụ thông tần xử lý cả hai vấn đề: nhiễu trục ngang của gia tốc kế khi không quay và sự trôi của con quay hồi chuyển.
\paragraph{}Ưu điểm:
\begin{itemize}
	\item Ước tính góc nhạy và chính xác, không bị ảnh hưởng nhiều với gia tốc ngang hoặc để con quay hồi chuyển bị trôi
	\item Không đòi hỏi chuyên sâu đối với bộ vi xử lý. Nó chỉ yêu cầu một lượng nhỏ các phép toán liên quan đến dấu chấm động. Có thể dễ dàng được chạy trong vòng lặp điều khiển bằng hoặc lớn hơn 100 Hz
	\item Trực quan và dễ giải thích về mặt lý thuyết hơn so với bộ lọc Kalman
\end{itemize}
\paragraph{}Nhược điểm:
\begin{itemize}
	\item Bản chất vẫn là từ các bộ lọc đơn giản như bộ lọc thông thấp và thông cao, tại các vùng nối tần số thường bị lệch pha rõ rệt
\end{itemize}
\paragraph{}Nói tóm lại, các bộ lọc thông thường là một kỹ thuật dùng phần cứng hoặc phần mềm nhằm giữ lạic các tín hiệu trong một khoảng thông dải tần số nào đó và loại bỏ tín hiệu ở các dải tần số còn lại. Đối với việc xây dựng bộ lọc bằng phần cứng tuy ra đời trước các bộ lọc phần mềm nhưng việc hiệu chỉnh đặc tính, thay đổi các tham số của bộ lọc phức tạp hơn rất nhiều so với sử dụng giải thuật xử lý tín hiệu số. Trong các bộ lọc, nếu tồn tại các tín hiệu nhiễu trong dải thông tần thì kết quả tín hiệu trở nên kém đi rất nhiều để có thể xử lý và điều khiển hệ thống một cách ổn định. Điều này tỏ ra càng thực tế đối với các bộ lọc phần cứng, vốn rất dễ bị nhiễu bởi các tín hiệu điện trong lúc hoạt động do sự kém chính xác của các linh kiện và sự bất ổn định của dòng điện ngõ vào. Trong nội dung đề tài này, bộ lọc được sử dụng là bộ lọc bổ phụ thông tần Complementary. 
\section{Mô hình lý thuyết động cơ DC}
\begin{wrapfigure}{r}[0pt]{0.4\linewidth}
	\begin{center}
		\includegraphics[scale=.75]{image/hinh32.png}
		\caption{Mô hình khung dây động cơ đặt trong từ trường}
	\end{center}
\end{wrapfigure}
\paragraph{}Động cơ DC có thể được mô hình hoá một cách đơn giản như hình trong hình 2.22. 
\paragraph{}Động cơ gồm một cuộn dây có N vòng dây quấn quanh một khung hình chữ nhật cạnh a, b, khung này có thể quay xung quanh một trục $\Delta$. Vị trí của nó được xác định bởi góc quay $\theta$. Cuộn dây có điện trở tổng cộng R và độ tự cảm L. 
\paragraph{}Hệ chuyển động có mô-men quán tính J với trục $\Delta$. Một nam châm vĩnh cửu tạo ra một từ trường B xuyên qua khung dây. 
\paragraph{}Hệ cơ S tác dụng lên trục một ngẫu lực cản, -$\Gamma$, được tính theo công thức: 
\begin{equation}
\Gamma = c\dot{\theta}
\end{equation}
với c là hằng số giảm chấn của kết cấu cơ khí (bộ truyền động ở trục động cơ). 
\newline
\begin{figure}[H]
	\begin{center}
		\includegraphics[scale=0.7]{image/hinh33.png}
		\caption{Mô hình điện cơ}
	\end{center}
\end{figure}
\paragraph{}Khi ta đặt điện áp một chiều E vào hai đầu chổi điện H và K, thì trong cuộn dây có dòng điện I$_{L}$. Theo hiện tượng cảm ứng điện từ, khung dây sẽ chịu tác dụng của lực điện trường (lực Laplace) và quay quanh trục $\Delta$. Ở đây H và K là hai vành bán khuyên, do đó cứ mỗi nửa vòng sẽ lại chuyển mạch một lần. Điều này đảm bảo cho khung luôn quay theo một chiều nhất định với vận tốc góc $\omega = \dot{\theta}$.
\paragraph{}E(t), I(t) là tổng điện áp cảm ứng và dòng dẫn trong động cơ. C(t), $\theta$'(t) là ngẫu lực điện từ và vận tốc xoay của rotor. Ta có công thức:
\begin{align}
E(t) = K\theta'(t)
\\
C(t) = KI(t)
\end{align}
với K là hằng số phụ thuộc và đặc tính của động cơ, loại cuộn dây, giá trị từ trường, thường là hằng số với nam châm vĩnh cửu. 
\paragraph{}Cuộn đây dẫn được mô hình hoá bởi điện trở R và độ dẫn L, được áp vào điện áp U(t), có phương trình:
\begin{equation}
U(t) = rI(t) + L\dfrac{d}{dt}I(t) + E(t)
\end{equation}
\paragraph{}Trong mối quan hệ ở trên, sự phụ thuộc vận tốc góc $\theta$'(t) theo C(t) là đặc tính cơ học của quán tính vật rắn và ma sát xoắn. Với J là mô-men quán tính, F$_{s}$(t) là ngẫu lực xoay và f$_{v}$ là hệ số ma sát quay, ta có phương trình cơ học:
\begin{equation}
J\dfrac{d}{dt}\theta(t) = C(t) - F_{s}(t) - f_{v}\theta_{v}\theta(t)
\end{equation}